随着我国跨区域人员流动日益频繁和公路货运需求持续增长，高速公路运行管理正由单纯保障通行逐步转向面向安全、效率与韧性的精细化治理。由于高速场景下车辆运行速度高、交通流密度大且扰动传播快，亟需构建高速路段持续感知与精细管理能力。智能交通系统（Intelligent Transportation Systems, ITS）的发展使交通状态实时感知与服务支撑成为可能，而能够连续描述车辆时空运动过程的轨迹数据已成为交通分析、预测与控制的重要基础。如何将路侧传感器输出的离散检测结果转化为连续、可靠的车辆轨迹，成为高速公路车辆目标跟踪研究需要解决的关键问题。

地磁传感器具有成本较低、体积较小、部署灵活、受光照和气象条件影响较小等优点，适用于低成本规模化部署。然而，在实际应用中仍面临两类问题：一类是通信链路波动、噪声干扰、邻道干扰及双车并行等因素导致的量测乱序、短时漏检与误检，降低在线跟踪准确度；另一类是传感器长期运行后检测性能衰减引起的连续漏检，易造成连续漏检条件下的轨迹碎片化，影响轨迹连续性与身份一致性。针对上述问题，本文通过在双车道道路两侧布设地磁传感器采集车辆信息，研究基于多地磁传感器的双车道车辆目标跟踪方法。本文主要研究内容如下：

针对双车道多地磁传感器车辆目标跟踪系统中的量测到达时序错乱、短时漏检、噪声干扰以及复杂并行车况，本文提出一种基于地磁传感器状态自适应与场景概率判别的在线多车辆目标跟踪方法。该方法首先利用时间窗口缓存对量测到达时序进行矫正，并通过固定滞后窗口接纳迟到到达的量测，结合局部回放机制提高轨迹信息的完整性；其次，在车辆运动模型与地磁量测模型基础上，利用卡尔曼滤波递推框架构建车辆状态估计方法，实现车辆位置与速度的在线预测和更新；随后，对跟踪波门及关联似然进行建模，并结合地磁传感器检测概率、杂波强度后验估计和地磁扰动峰值一致性信息计算关联权重，实现杂波环境下车辆轨迹与多地磁传感器量测之间的稳健关联；最后，构建邻道干扰与双车并行三状态场景判别模型，并结合车辆状态递推更新，对车辆行驶车道进行在线判别与修正。实验结果表明，所提方法能够在地磁传感器检测数据乱序、漏检及噪声干扰条件下，实现对车辆位置、速度和行驶车道的较准确估计，并有效区分双车并行与邻道干扰两类易混淆场景。

针对系统长期运行后部分地磁传感器检测能力下降、连续漏检增多所引发的长缺测轨迹碎片化问题，本文进一步提出连续漏检场景下的多地磁传感器车辆轨迹实时聚类算法。该方法以上游在线跟踪输出的轨迹片段及未关联的单点量测作为输入，构建适用于在线增量处理的簇图结构，并围绕簇初始化、聚类相似度评分、在线归簇以及簇状态更新与生命周期维护等关键环节，形成片段级轨迹重构流程。实验结果表明，所提方法能够在连续漏检条件下有效恢复断裂轨迹片段，降低错误拼接和身份切换，提升轨迹连续性与身份一致性。